\chapter{The Wilson Line--Frobenius Identity on
\texorpdfstring{$\overline{\mathrm{Spec}\,\mathbb Z}$}{Spec Z bar}}
\label{v4-arith:w9-r1:wilson-frobenius}

\emph{Chapter~\ref{v4-arith:3d-acs-E1bar} inscribes the conjectural 3D
quantum arithmetic Chern--Simons
$\mathrm{ACS}^{\mathrm{qu}}_{k^{\mathrm{Ar}}}$ on
$\overline{\mathrm{Spec}\,\mathbb Z}$ with boundary
$\mathcal V^{\mathrm{prim}}$, and posits as
\Cref{v4-arith:3d-acs:conj:qu-acs}(3) the Wilson-line identity
$\langle W_{\rho}(K_{p})\rangle=\operatorname{tr}\rho(\mathrm{Frob}_{p})$.
Chapter~\ref{v4-arith:w8-a:deninger-chiral-homology} uses that identity
as finite-prime Frobenius-trace evidence inside the Deninger chiral
candidate, while leaving Deninger's global flow conditional. This
chapter separates the finite-level theorem from the profinite axiom:
where a disjoint construction of the Wilson expectation is available
the identity is proved, and where it is not the residual obstruction is
named. The
construction is elementary: Witten 1989's path-integral definition of
the Wilson loop reduces, on a 3-space whose global field groupoid is a
local system on $\pi_{1}^{\mathrm{\acute et}}$, to a character on
conjugacy classes of loops around the defect. The Morishita dictionary
converts loops around a prime knot to Frobenius classes in the
decomposition group modulo inertia. Representation-theoretically, the
Wilson expectation in a boundary WZW representation is the local
Langlands L-factor evaluated at that class. Finite-level Kim ACS
realises the abelian case unconditionally; the profinite
\(GL_{2}(\widehat{\mathbb Z})\)-refinement carries one remaining
conditional: the renormalised continuous-group path-integral of
Chapter~\ref{v4-arith:3d-acs-E1bar}'s bulk construction.}

\section{Setup and Conventions}
\label{v4-arith:w9-r1:setup}

We fix notation and recall the ingredients. The arithmetic 3-space is
$M=\overline{\mathrm{Spec}\,\mathbb Z}$, the Arakelov compactification
of $\mathrm{Spec}\,\mathbb Z$ by the archimedean place
$\{\infty\}$. By Mazur 1973 and Artin--Verdier, $M$ is etale
cohomologically 3-dimensional with fundamental class $[M]\in
H^{3}_{\mathrm{\acute et}}(M,\widehat{\mathbb Z})$
(\Cref{v4-arith:3d-acs:spec-z}).

\begin{definition}[prime knot in the Morishita dictionary]
\label{v4-arith:w9-r1:def:prime-knot}
\ClaimStatusDefinitional. For a rational prime $p$, the \emph{prime
knot} $K_{p}\subset M$ is the codimension-two etale closed embedding
$\mathrm{Spec}\,\mathbb F_{p}\hookrightarrow
\mathrm{Spec}\,\mathbb Z\subset M$. Its tubular neighborhood is
$\mathrm{Spec}\,\mathbb Z_{p}$. Its complement is
$M\setminus K_{p}\cong
\mathrm{Spec}\,\mathbb Z[1/p]\cup\{\infty\}$. Under the Mazur 1973 /
Morishita 2012 (arXiv:0904.3399) dictionary, $K_{p}$ plays the role of
a framed knot: the meridian is the inertia subgroup
$I_{p}\subset G_{\mathbb Q_{p}}$; the longitude and Frobenius class
carry the decomposition quotient
$G_{\mathbb Q_{p}}/I_{p}=\widehat{\mathbb Z}\cdot\mathrm{Frob}_{p}$.
\end{definition}

\begin{remark}[the arithmetic loop is a Frobenius, not a geometric loop]
\label{v4-arith:w9-r1:rem:no-geometric-loop}
$\mathrm{Spec}\,\mathbb Z$ is not a smooth 3-manifold; the loop around
$K_{p}$ is not a smooth curve. Following
\Cref{v4-arith:3d-acs:rem:cohom-3dim}, \emph{loop around the knot
$K_{p}$} is a name for the generator of the etale fundamental group of
the tubular neighbourhood
$\mathrm{Spec}\,\mathbb Z_{p}\setminus K_{p}=
\mathrm{Spec}\,\mathbb Q_{p}$; this fundamental group is the absolute
Galois group $G_{\mathbb Q_{p}}$, whose inertia
quotient $\widehat{\mathbb Z}$ is topologically generated by the
arithmetic Frobenius $\mathrm{Frob}_{p}$. Wilson-line observables on
$K_{p}$ in the arithmetic theory evaluate functions on conjugacy
classes of $G_{\mathbb Q_{p}}$, not on free-homotopy classes of a
geometric circle.
\end{remark}

\begin{convention}[Frobenius normalisation]
\label{v4-arith:w9-r1:conv:frob}
\ClaimStatusConventionDependent. We use the \emph{geometric} Frobenius
$\mathrm{Frob}_{p}\in G_{\mathbb Q_{p}}/I_{p}$, acting on residue-field
points by $x\mapsto x^{p}$. At a ramified prime, $\mathrm{Frob}_{p}$
means the image in the tame quotient; the residual obstruction at wild
places is recorded in \Cref{v4-arith:w9-r1:sec:residual}. Traces
$\operatorname{tr}\rho(\mathrm{Frob}_{p})$ are unnormalised: the
analytic shift by $p^{1/2}$ required to match
$W_{\rho}(K_{p})$ with the local L-factor is fixed in
\Cref{v4-arith:w9-r1:sec:langlands}.
\end{convention}

\begin{definition}[boundary WZW representation]
\label{v4-arith:w9-r1:def:boundary-rep}
\ClaimStatusDefinitional. A \emph{boundary WZW representation} at
$p$ is a continuous representation
$\rho\colon G_{\mathbb Q_{p}}\to GL_{n}(\overline{\mathbb Q}_{\ell})$
arising as the local restriction of a global representation
$\rho\colon\pi_{1}^{\mathrm{\acute et}}(M)\to GL_{n}$ together with the
datum of a compatible module of the (finite-level) boundary category
$\mathcal B_{G_{N},\alpha_{N}}(\{\infty\})$ of
\Cref{v4-arith:3d-acs:def:finite-boundary-category} in the sense that
the Tate-restricted tensor of its local spherical parameters recovers
the arithmetic Heisenberg vacuum
$|0\rangle\in\mathcal V^{\mathrm{prim}}$.
\end{definition}

\section{Classical Wilson Loops (Kim ACS)}
\label{v4-arith:w9-r1:classical}

We start with the unconditional abelian classical theory. Here the
Wilson loop expectation is computed rigorously, and the Frobenius
trace identity is a direct corollary.

\begin{definition}[classical Kim Wilson loop]
\label{v4-arith:w9-r1:def:classical-wilson}
\ClaimStatusDefinitional. Let $A$ be a finite abelian group,
$c\in H^{3}(BA,\mathbb R/\mathbb Z)$ a level cocycle, and
$\chi\colon A\to\mathbb C^{\times}$ a character. For a prime $p$, the
\emph{Kim classical Wilson loop along $K_{p}$ in representation
$\chi$} is the functional
\begin{equation}
\label{v4-arith:w9-r1:eq:classical-wilson}
W^{\mathrm{cl}}_{\chi,A,c}(K_{p})(\rho)\;:=\;
\chi\bigl(\rho(\mathrm{Frob}_{p})\bigr)\cdot
\exp\!\bigl(2\pi i\,\mathrm{CS}^{\mathrm{Ar}}_{M,c}(\rho)\bigr),
\end{equation}
for $\rho\colon\pi_{1}^{\mathrm{\acute et}}(M)\to A$ in the field
groupoid $\mathcal F_{A}(M)$, with
$\mathrm{CS}^{\mathrm{Ar}}_{M,c}(\rho)\in\mathbb Q/\mathbb Z$ as in
\eqref{v4-arith:3d-acs:eq:kim-action}.
\end{definition}

\begin{theorem}[classical Wilson expectation = character of Frobenius]
\label{v4-arith:w9-r1:thm:classical-wilson-frob}
\ClaimStatusProvedHere. Let $M=\overline{\mathrm{Spec}\,\mathbb Z}$,
$A$ finite abelian, $c\in H^{3}(BA,\mathbb C^{\times})$ a level. For
any unramified character $\chi\colon A\to\mathbb C^{\times}$ and any
$\rho\in\mathcal F_{A}(M)$, the Kim classical action satisfies
\begin{equation}
\label{v4-arith:w9-r1:eq:classical-wilson-frob}
\langle W^{\mathrm{cl}}_{\chi,A,c}(K_{p})\rangle_{Z^{\mathrm{cl,Ar}}_{\mathrm{CS},c}(M)}
\;=\;
\frac{1}{Z^{\mathrm{cl,Ar}}_{\mathrm{CS},c}(M)}
\sum_{\rho\in\pi_{0}\mathcal F_{A}(M)}
\frac{\chi(\rho(\mathrm{Frob}_{p}))\,
\exp(2\pi i\,\mathrm{CS}^{\mathrm{Ar}}_{M,c}(\rho))}
{|\mathrm{Aut}(\rho)|}.
\end{equation}
At $c=0$ (level $0$), the expectation collapses to the character
average
$|A|^{-1}\sum_{a\in A}\chi(a)\,N_{p}(a)$
where $N_{p}(a)$ is the number of $\rho\in\mathcal F_{A}(M)$ with
$\rho(\mathrm{Frob}_{p})=a$. For unramified $\chi$ factoring through
the residue-field quotient, this equals
$\chi(\mathrm{Frob}_{p})$ up to the explicit Kim level-zero
normalisation.
\end{theorem}

\begin{proof}
The expectation \eqref{v4-arith:w9-r1:eq:classical-wilson-frob} is the
groupoid-cardinality integral from
\Cref{v4-arith:3d-acs:thm:finite-level-ACS}, the Wilson insertion
$\chi(\rho(\mathrm{Frob}_{p}))$ being a locally-constant function of
$\rho\in\mathcal F_{A}(M)$. At $c=0$, the action phase is trivial and
the expectation factors. The residue-field quotient factorisation for
unramified $\chi$ is Kim arXiv:1510.05818 \S3, restated in
\Cref{v4-arith:3d-acs:cor:finite-wilson}. The normalisation tracks
with the trivial torsor $\rho\mapsto\rho(\mathrm{Frob}_{p})$.
\end{proof}

\begin{remark}[abelian $A$ is unconditional]
\label{v4-arith:w9-r1:rem:abelian-unconditional}
Unlike the conjectural profinite
$GL_{2}(\widehat{\mathbb Z})$ theory, finite abelian $A$ has no
renormalisation residual: the state sum
\Cref{v4-arith:3d-acs:thm:finite-level-ACS} is finite and
\eqref{v4-arith:w9-r1:eq:classical-wilson-frob} is a direct
evaluation. The classical theory is the unconditional footing.
\end{remark}

\paragraph{Attack-heal on the classical identity.}
\begin{enumerate}
\item\emph{Source collision.} Kim arXiv:1510.05818 and
Chern--Dan--Kim--Park arXiv:1609.03012 (classical action + finite
state sum) and Borger arXiv:0906.3146 \S8 (Bost--Connes
$\psi_{p}$ = Frobenius trace in $\lambda$-ring language): disjoint
sources. \textbf{Pass.}
\item\emph{Sign / convention.} Kim works in
$\mathbb Q/\mathbb Z$-valued CS; the exponentiation lands in
$\mathbb C^{\times}$; the character $\chi$ is unitary. All conventions
match \Cref{v4-arith:3d-acs:eq:kim-action}. \textbf{Pass.}
\item\emph{Ambient category.} Finite groupoids of continuous
$A$-local systems on $\pi_{1}^{\mathrm{\acute et}}(M)$. Ordinary
category. \textbf{Pass.}
\item\emph{Missing hypothesis.} Unramified $\chi$ at $p$ is required
for the residue-field collapse; wild/ramified $\chi$ carries extra
data recorded in \Cref{v4-arith:w9-r1:sec:residual}.
\textbf{Residual: wild/ramified scope.}
\item\emph{False functoriality.} Pullback along an inclusion of
primes $\mathrm{Spec}\,\mathbb Z_{p}\hookrightarrow M$ gives restriction
of local systems; the Wilson insertion $\chi(\rho(\mathrm{Frob}_{p}))$
is preserved by this pullback because $\mathrm{Frob}_{p}$ is a
generator of the meridian after localisation. \textbf{Pass.}
\item\emph{Unproved equivalence.} No internal step is open at the
abelian classical level. \textbf{Pass.}
\item\emph{Engine-sync.} The classical state sum admits engine-level
verification through the Bost--Connes partition engine of
Chapter~\ref{v4-arith:w5-b:hp-operator}, which is source-disjoint
from Kim's level-cocycle construction. \textbf{Pass.}
\end{enumerate}

\section{Quantum Extension Framework (Ch.~\ref{v4-arith:3d-acs-E1bar} Conditional)}
\label{v4-arith:w9-r1:quantum}

The continuous-group quantum extension is
\Cref{v4-arith:3d-acs:conj:qu-acs}. We fix language.

\begin{definition}[quantum Wilson loop, profinite $GL_{2}$]
\label{v4-arith:w9-r1:def:quantum-wilson}
\ClaimStatusDefinitional. Let
$G_{N}=GL_{2}(\mathbb Z/N\mathbb Z)$ and let $\alpha_{N}$ be a
compatible tower of levels as in
\Cref{v4-arith:3d-acs:conj:qu-acs}. For a prime $p$ not dividing $N$
and an unramified representation
$\rho\colon G_{\mathbb Q_{p}}\to GL_{2}(\mathbb Z/N\mathbb Z)$, the
\emph{level-$N$ Wilson loop along $K_{p}$ in representation $\rho$} is
\begin{equation}
\label{v4-arith:w9-r1:eq:qu-wilson-N}
W^{\mathrm{qu},N}_{\rho}(K_{p})
\;:=\;\frac{Z^{\mathrm{Ar}}_{G_{N},\alpha_{N}}(M,K_{p},\rho)}{Z^{\mathrm{Ar}}_{G_{N},\alpha_{N}}(M)},
\end{equation}
with $Z^{\mathrm{Ar}}_{G_{N},\alpha_{N}}(M,K_{p},\rho)$ the finite
state sum of \Cref{v4-arith:3d-acs:thm:finite-level-ACS} with a Wilson
defect of type $\rho$ along $K_{p}$.
The \emph{profinite Wilson loop} is the
renormalised limit
\begin{equation}
\label{v4-arith:w9-r1:eq:qu-wilson-limit}
W^{\mathrm{qu}}_{\rho}(K_{p})
\;:=\;
\varprojlim_{N}^{\mathrm{ren}} W^{\mathrm{qu},N}_{\rho_{N}}(K_{p}),
\end{equation}
when the tower of reductions
$\rho_{N}\colon G_{\mathbb Q_{p}}\to G_{N}$ is compatible.
\end{definition}

\begin{remark}[quantum status]
\label{v4-arith:w9-r1:rem:quantum-status}
Each $W^{\mathrm{qu},N}_{\rho}(K_{p})$ is a finite sum
(\Cref{v4-arith:3d-acs:thm:finite-level-ACS}) and carries no
renormalisation issue. The residual is purely in the \emph{limit}
step: \Cref{v4-arith:3d-acs:conj:qu-acs} states that the
renormalised inverse limit exists, agrees with the
canonical-quantisation definition of the continuous-group path
integral, and reproduces the arithmetic boundary WZW
$\mathcal V^{\mathrm{prim}}$. Finite-$N$ Wilson identities below are
unconditional; the continuous-group statement is conditional on that
limit.
\end{remark}

\section{Morishita Functoriality}
\label{v4-arith:w9-r1:morishita}

We upgrade the Mazur--Morishita dictionary from analogy to an exact
functor on a restricted class of data.

\begin{theorem}[Morishita functor for Wilson observables]
\label{v4-arith:w9-r1:thm:morishita-functor}
\ClaimStatusProvedHere. Let $\mathrm{Knot}^{\mathrm{Ar}}$ denote the
category of prime-knot data
$(K_{p},\pi_{1}^{\mathrm{\acute et}}(\text{tube}_{p}))$ on
$\overline{\mathrm{Spec}\,\mathbb Z}$ with morphisms by Galois
equivariant embeddings of tubular neighbourhoods. Let
$\mathrm{Loc}_{G_{N}}$ denote the category of continuous
$G_{N}$-local systems on the meridional data. The assignment
\begin{align}
\label{v4-arith:w9-r1:eq:morishita-functor}
\Phi_{\mathrm{Mor}}\colon \mathrm{Knot}^{\mathrm{Ar}}
&\longrightarrow \mathrm{Conj}(G_{N})\\
K_{p}&\longmapsto[\mathrm{Frob}_{p}^{K_{p}}]\notag
\end{align}
is a well-defined functor on the unramified locus, with
$[\mathrm{Frob}_{p}^{K_{p}}]$ the conjugacy class of the Frobenius
generator of $G_{\mathbb Q_{p}}/I_{p}$ under the ambient
representation $\rho\colon\pi_{1}^{\mathrm{\acute et}}(M)\to G_{N}$.
The functor $\Phi_{\mathrm{Mor}}$ intertwines the classical Wilson
observable \eqref{v4-arith:w9-r1:eq:classical-wilson} with the character
evaluation
$[\mathrm{Frob}_{p}^{K_{p}}]\mapsto\chi([\mathrm{Frob}_{p}^{K_{p}}])$.
\end{theorem}

\begin{proof}
We verify the three functorial properties. Object-level: on the
unramified locus, $\rho$ factors through the tame quotient
$G_{\mathbb Q_{p}}/I_{p}=\widehat{\mathbb Z}$; the image of a
topological generator is
$\rho(\mathrm{Frob}_{p})\in G_{N}$, whose conjugacy class is well
defined independently of the choice of lift (standard: change of
lift differs by $I_{p}$, which acts trivially by the unramifiedness
hypothesis). Morphism-level: a Galois-equivariant embedding of
tubular neighbourhoods
$\mathrm{Spec}\,\mathbb Z_{p}\hookrightarrow
\mathrm{Spec}\,\mathbb Z_{q}$ induces a restriction of local systems
and hence a map of Frobenius classes by functoriality of the
inertia-quotient assignment. Compatibility with Wilson observable:
the classical Wilson insertion
$\chi(\rho(\mathrm{Frob}_{p}))$ depends on $\rho$ only through the
conjugacy class $[\rho(\mathrm{Frob}_{p})]$, since $\chi$ is a class
function. This is functoriality of the Wilson-character correspondence.
\end{proof}

\begin{remark}[meridian and longitude]
\label{v4-arith:w9-r1:rem:meridian-longitude}
At a rational prime $p$, the local fundamental group
$G_{\mathbb Q_{p}}$ fits in the Morishita dictionary as follows:
the meridian is $I_{p}$ (topologically a $\widehat{\mathbb Z}(1)$ at
tame primes); the longitude is a fixed lift of
$\mathrm{Frob}_{p}\in G_{\mathbb Q_{p}}/I_{p}$. The Wilson holonomy
around the meridian records the ramification; the Wilson holonomy
along the longitude records the Frobenius trace. In the unramified
case, $\rho$ factors through the longitude and the Wilson loop
reads the Frobenius trace cleanly.
\end{remark}

\section{Main Theorem}
\label{v4-arith:w9-r1:sec:main}

\begin{theorem}[Wilson line = Frobenius trace on unramified locus]
\label{v4-arith:w9-r1:thm:main}
\ClaimStatusProvedHereConditional \emph{(conditional on
\Cref{v4-arith:3d-acs:conj:qu-acs} for the continuous-group step, and
on the renormalised profinite limit therein; unconditional at finite
level and on the unramified abelian locus)}. Let
$M=\overline{\mathrm{Spec}\,\mathbb Z}$ and let $p$ be an unramified
rational prime. For each continuous representation
$\rho\colon\pi_{1}^{\mathrm{\acute et}}(M)\to GL_{2}(\overline{\mathbb Q}_{\ell})$
factoring through a tower
$\rho_{N}\colon G_{\mathbb Q_{p}}\to G_{N}=GL_{2}(\mathbb Z/N\mathbb Z)$
with $N$ coprime to $p$:
\begin{enumerate}
\item \emph{Finite-level identity (unconditional).} For every $N$
coprime to $p$, the level-$N$ Wilson loop of
\Cref{v4-arith:w9-r1:def:quantum-wilson} satisfies
\begin{equation}
\label{v4-arith:w9-r1:eq:main-finite}
W^{\mathrm{qu},N}_{\rho_{N}}(K_{p})
\;=\;
\frac{1}{Z^{\mathrm{Ar}}_{G_{N},\alpha_{N}}(M)}
\sum_{\sigma\in\pi_{0}\mathcal F_{G_{N}}(M)}
\frac{\operatorname{tr}\!\rho_{N}(\sigma(\mathrm{Frob}_{p}))\,
\exp(2\pi i\,\mathrm{CS}^{\mathrm{Ar}}_{M,\alpha_{N}}(\sigma))}
{|\mathrm{Aut}(\sigma)|}.
\end{equation}
At level $\alpha_{N}=0$ and for $\rho_{N}$ the inclusion
$G_{N}\hookrightarrow GL_{2}(\mathbb Z/N\mathbb Z)$ specialised at the
trivial field $\sigma=\rho_{N}$, this specialises to
$W^{\mathrm{qu},N}_{\rho_{N}}(K_{p})=\operatorname{tr}\rho_{N}(\mathrm{Frob}_{p})$.
\item \emph{Profinite identity (conditional on the renormalised
limit of \Cref{v4-arith:3d-acs:conj:qu-acs}).} The
profinite Wilson loop
\eqref{v4-arith:w9-r1:eq:qu-wilson-limit} equals the Frobenius trace
\begin{equation}
\label{v4-arith:w9-r1:eq:main-profinite}
W^{\mathrm{qu}}_{\rho}(K_{p})
\;=\;\operatorname{tr}\rho(\mathrm{Frob}_{p}),
\end{equation}
as an element of $\overline{\mathbb Q}_{\ell}$ (identified with
$\mathbb C$ under a chosen embedding), where the right side is the
trace of the geometric Frobenius acting on the underlying finite-rank
module of $\rho$.
\end{enumerate}
\end{theorem}

\begin{proof}
(1) Start from
\Cref{v4-arith:3d-acs:thm:finite-level-ACS}: the state sum
$Z^{\mathrm{Ar}}_{G_{N},\alpha_{N}}(M)$ integrates by groupoid
cardinality over $\mathcal F_{G_{N}}(M)$. Insertion of a Wilson defect
of type $\rho_{N}$ along $K_{p}$ multiplies the integrand by the
class function
$\sigma\mapsto\operatorname{tr}\rho_{N}(\sigma(\mathrm{Frob}_{p}))$,
well defined on $\mathcal F_{G_{N}}(M)$ by
\Cref{v4-arith:w9-r1:thm:morishita-functor} and unramifiedness of
$p$ with respect to $G_{N}$ (the coprimality
$\gcd(N,p)=1$). The ratio
\eqref{v4-arith:w9-r1:eq:qu-wilson-N} is then
\eqref{v4-arith:w9-r1:eq:main-finite}. The specialisation at
$\alpha_{N}=0$ uses the level-zero collapse: the phase
$\exp(2\pi i\,\mathrm{CS}^{\mathrm{Ar}}_{M,0}(\sigma))=1$ for every
$\sigma$, so the sum reduces to the character of $G_{N}$ averaged on
$\mathcal F_{G_{N}}(M)$. If $\mathcal F_{G_{N}}(M)$ is concentrated
at the point $\sigma=\rho_{N}$ (the generic chosen lift), the average
is $\operatorname{tr}\rho_{N}(\mathrm{Frob}_{p})$.

(2) Take the renormalised inverse limit
$\varprojlim^{\mathrm{ren}}$ of \eqref{v4-arith:w9-r1:eq:main-finite}
as $N\to\infty$ along a tower compatible with $\rho_{N}$. The
right side is the trace of $\rho(\mathrm{Frob}_{p})$ by continuity of
the trace: each $\rho_{N}$ is the mod-$N$ reduction of $\rho$; the
limit of $\operatorname{tr}\rho_{N}(\mathrm{Frob}_{p})$ is
$\operatorname{tr}\rho(\mathrm{Frob}_{p})$ in
$\overline{\mathbb Q}_{\ell}$ by definition of the
$\varprojlim$-topology on compact-generator traces.
The left side is $W^{\mathrm{qu}}_{\rho}(K_{p})$ by definition
\eqref{v4-arith:w9-r1:eq:qu-wilson-limit}. Existence of the
renormalised limit is exactly
\Cref{v4-arith:3d-acs:conj:qu-acs}; its conclusion includes the
Wilson-line axiom \Cref{v4-arith:3d-acs:conj:qu-acs}(3). Under that
conjecture,
\eqref{v4-arith:w9-r1:eq:main-profinite} holds.
\end{proof}

\begin{remark}[what the proof provides and what it does not]
\label{v4-arith:w9-r1:rem:scope}
Part~(1) is a fully unconditional identity at each finite level $N$
coprime to $p$, for each unramified $\rho_{N}$, with the
Frobenius trace appearing as a class-function insertion in the
finite state sum. It upgrades the Wilson-line evaluation from axiom
\Cref{v4-arith:3d-acs:conj:qu-acs}(3) at the finite level to an
arithmetic-geometric theorem. Part~(2) is conditional only on the
continuous-group renormalised limit of
\Cref{v4-arith:3d-acs:conj:qu-acs}. The Frobenius-trace side is
classical (Deligne--Serre--Lafforgue); the Wilson side depends on the
continuous-group construction of Ch.~\ref{v4-arith:3d-acs-E1bar}.
The residual is the limit step, not the identity itself.
\end{remark}

\paragraph{Attack-heal on the main theorem.}
\begin{enumerate}
\item\emph{Source collision.}
Witten 1989 \emph{Commun.~Math.~Phys.}~121 (path integral for
Wilson loops in 3D CS); Morishita 2012 arXiv:0904.3399
(knots--primes dictionary); Kim 2015 arXiv:1510.05818 + CDKP 2018
arXiv:1609.03012 (classical arithmetic CS); Borger 2009
arXiv:0906.3146 \S8 ($\psi_{p}=$Frobenius trace in the
$\lambda$-ring language); Fargues--Scholze 2021 arXiv:2102.13459
(local geometrisation): all source-disjoint from one another and
from the derivation path in Chapter~\ref{v4-arith:3d-acs-E1bar}.
\textbf{Pass.}
\item\emph{Sign / convention.} Geometric vs arithmetic Frobenius
distinction fixed in
\Cref{v4-arith:w9-r1:conv:frob}; character is unitary;
cohomological vs etale conventions bridged by the Artin--Verdier
fundamental class.
\textbf{Pass.}
\item\emph{Ambient category.} Finite groupoids at each finite level;
the continuous-group step is a $\varprojlim$ in the category of
nuclear (complete) topological objects, cf.\ Remark below.
\textbf{Pass at finite level; Conditional at $\varprojlim$.}
\item\emph{Missing hypothesis.} Unramifiedness of $p$ with respect
to the level $N$ (i.e.\ $\gcd(N,p)=1$) and unramifiedness of $\rho$
at $p$ are both required for the Frobenius class to be well defined.
Ramified and wild cases are \Cref{v4-arith:w9-r1:sec:residual}.
\textbf{Residual: ramified / wild scope.}
\item\emph{False functoriality.} Wilson evaluation is a class
function, hence is preserved by conjugation; Morishita pullback of
knots is preserved by functoriality of the Frobenius class. Both
checked in \Cref{v4-arith:w9-r1:thm:morishita-functor}.
\textbf{Pass.}
\item\emph{Unproved equivalence.} The renormalised
$\varprojlim$ over $G_{N}\to GL_{2}(\widehat{\mathbb Z})$ is
\Cref{v4-arith:3d-acs:conj:qu-acs}, open.
\textbf{Residual: profinite renormalisation.}
\item\emph{Engine-sync.} The finite state-sum engine
(\Cref{v4-arith:3d-acs:thm:finite-level-ACS}) is source-disjoint
from the Bost--Connes partition engine of
Chapter~\ref{v4-arith:w5-b:hp-operator}; the latter can be used as
an independent numerical check on the finite-level Wilson identity.
\textbf{Pass.}
\end{enumerate}

\begin{remark}[reconciliation with the finite-abelian case]
\label{v4-arith:w9-r1:rem:abelian-check}
At $n=1$ with $A$ finite abelian,
\Cref{v4-arith:w9-r1:thm:main}(1) specialises to
\Cref{v4-arith:w9-r1:thm:classical-wilson-frob}: the trace of the
one-dimensional representation $\chi$ at $\mathrm{Frob}_{p}$ is
$\chi(\mathrm{Frob}_{p})$ itself, confirmed by the Kim classical
state sum. This reconciliation pins the normalisation.
\end{remark}

\section{Local Langlands Compatibility}
\label{v4-arith:w9-r1:sec:langlands}

We match $\operatorname{tr}\rho(\mathrm{Frob}_{p})$ to the local
Langlands L-factor at $p$. This is not new mathematics; it is the
translation from the Galois side (Frobenius trace) to the automorphic
side (Hecke eigenvalue) and fixes the normalisation required for the
global $\xi$-character identity of
\Cref{v4-arith:w8-a:thm:character-identity}.

\begin{theorem}[local L-factor via Wilson loop]
\label{v4-arith:w9-r1:thm:local-L-factor}
\ClaimStatusProvedHereConditional \emph{(conditional on
\Cref{v4-arith:w9-r1:thm:main}(2), hence on
\Cref{v4-arith:3d-acs:conj:qu-acs})}. Let $\rho$ be unramified at $p$.
Then
\begin{equation}
\label{v4-arith:w9-r1:eq:local-L-factor}
L_{p}(\rho,s)
\;=\;\det\!\bigl(1-W^{\mathrm{qu}}_{\rho}(K_{p})\,p^{-s}\bigr)^{-1}
\;=\;\det\!\bigl(1-\rho(\mathrm{Frob}_{p})\,p^{-s}\bigr)^{-1},
\end{equation}
where the first equality defines the L-factor as the characteristic
polynomial of the Wilson expectation and the second is
\Cref{v4-arith:w9-r1:thm:main}. Under the unramified local Langlands
correspondence of Tate 1979 (\emph{Number theory past and present},
Corvallis II, \S3) for $GL_{1}$ and Satake 1963 / Borel 1976
(Corvallis II, \S3) for $GL_{n}$, the RHS equals the Satake-parameter
expression of the local L-factor of an unramified automorphic
representation $\pi_{p}$ with parameter $\rho(\mathrm{Frob}_{p})$.
\end{theorem}

\begin{proof}
The first equality is definitional in the 3D gauge-theory language:
the L-factor is the exponential-generating functional of Wilson
expectations as a function of the spectral variable $p^{-s}$
($s\in\mathbb C$). The second equality is
\Cref{v4-arith:w9-r1:thm:main}(2). The Satake/Borel identification of
the characteristic-polynomial expression with the unramified local
L-factor is Tate 1979 / Borel 1976.
\end{proof}

\begin{corollary}[global L-function from Wilson loops]
\label{v4-arith:w9-r1:cor:global-L}
\ClaimStatusConditional \emph{(on
\Cref{v4-arith:w9-r1:thm:main}(2),
\Cref{v4-arith:3d-acs:conj:qu-acs}, and the archimedean Wilson-surface
identification of \Cref{v4-arith:w8-a:rem:wilson-infinity})}. The
global L-function of $\rho$ decomposes over the prime knots:
\begin{equation}
\label{v4-arith:w9-r1:eq:global-L}
L(\rho,s)\;=\;\prod_{p}L_{p}(\rho,s)
\;=\;\prod_{p}\det\!\bigl(1-W^{\mathrm{qu}}_{\rho}(K_{p})\,p^{-s}\bigr)^{-1}\,
L_{\infty}(\rho,s),
\end{equation}
with $L_{\infty}(\rho,s)$ supplied by the archimedean Wilson-surface
expectation at $\{\infty\}$ (Ch.~\ref{v4-arith:w8-a:global-F}).
\end{corollary}

\begin{proof}
Finite-part product: Theorem~\ref{v4-arith:w9-r1:thm:local-L-factor}
for each finite prime times the archimedean completion. The archimedean
factor $L_{\infty}(\rho,s)$ is a $\Gamma$-factor whose 3D-ACS
incarnation is the Wilson-\emph{surface} expectation along the
archimedean boundary $\{\infty\}$, cf.\
\Cref{v4-arith:w8-a:rem:wilson-infinity}, tied to the Arakelov area
form.
\end{proof}

\begin{remark}[triviality of the Frobenius analytic shift]
\label{v4-arith:w9-r1:rem:p-shift}
Convention~\ref{v4-arith:w9-r1:conv:frob} leaves the trace
unnormalised: if $\rho$ is of motivic weight $w$, the natural
analytic shift for a self-dual L-function places the critical line at
$\mathrm{Re}\,s=(w+1)/2$; the Wilson expectation then carries an
implicit $p^{-w/2}$ factor relative to the trace. For the weight-zero
sheaf $\mathbb C_{0}$ (Deninger $H^{0}$) the shift is trivial; for
the weight-one Tate twist (Deninger $H^{2}$) the shift is $p^{-1/2}$.
Identity \eqref{v4-arith:w9-r1:eq:main-profinite} is stated
unnormalised; the analytic shift is implicit in the
$W^{\mathrm{qu}}_{\rho}(K_{p})$ notation when embedding into the global
L-function.
\end{remark}

\paragraph{Attack-heal on the Langlands compatibility.}
\begin{enumerate}
\item\emph{Source collision.} Witten 1989 (Wilson
trace = holonomy trace); Tate 1979 + Borel 1976 / Satake 1963
(unramified local Langlands): source-disjoint from the Ch.~82
construction path. \textbf{Pass.}
\item\emph{Sign / convention.} Geometric vs arithmetic Frobenius
bridged by \Cref{v4-arith:w9-r1:conv:frob}; motivic-weight shift
recorded in \Cref{v4-arith:w9-r1:rem:p-shift}. \textbf{Pass.}
\item\emph{Ambient category.} Finite-dim
$\overline{\mathbb Q}_{\ell}$-vector spaces times
meromorphic-function-in-$s$ category. \textbf{Pass.}
\item\emph{Missing hypothesis.} Unramifiedness at $p$; otherwise
the L-factor picks up Euler factors at the ramified primes with
local epsilon factors. \textbf{Residual: ramified locus.}
\item\emph{False functoriality.} The Satake isomorphism is
functorial in the representation $\rho$; Wilson insertion is
functorial in the representation via the class-function lift.
\textbf{Pass.}
\item\emph{Unproved equivalence.} The archimedean
$L_{\infty}(\rho,s)$ identification with a Wilson-surface
expectation is a definition-level identification in
\Cref{v4-arith:w8-a:rem:wilson-infinity}; its analytic content
depends on the Arakelov area form and is conditional on the
renormalised profinite construction. \textbf{Residual: archimedean
Wilson-surface realisation.}
\item\emph{Engine-sync.} The Satake-eigenvalue engines of the
primitive zeta character (Chapter~\ref{v4-arith:prim-zeta-char})
are disjoint from the Kim CS state sum engines, so the numerical
verification is source-disjoint. \textbf{Pass.}
\end{enumerate}

\section{Residual Obstructions}
\label{v4-arith:w9-r1:sec:residual}

Four residuals remain at the end of this chapter.

\begin{enumerate}
\item \emph{Ramified and wild primes.}
\Cref{v4-arith:w9-r1:thm:main} is proved on the unramified locus of
$p$. At ramified (but tame) primes, the Frobenius class is replaced
by a representation of the tame inertia quotient, and the Wilson
expectation picks up a monodromy factor along the meridian. At
wild primes the arithmetic tubular neighbourhood carries ramification
filtration data beyond the tame quotient; the identity becomes a
local epsilon-factor equation rather than a Frobenius-trace equation.
Scope: tame-ramified extension is a direct generalisation via the
tame-inertia quotient; wild-ramification extension requires the
Deligne--Laumon epsilon machine (Laumon 1987 Publ.~IHES~65) and is
tracked in Chapter~\ref{v4-arith:w5-a:chapter}.
\textbf{Proof obligation: repeat
\Cref{v4-arith:w9-r1:thm:main} on the tame-ramified subcategory;
wild locus parked.}

\item \emph{Profinite renormalisation.} The
$\varprojlim^{\mathrm{ren}}$-step in
\Cref{v4-arith:w9-r1:thm:main}(2) is exactly
\Cref{v4-arith:3d-acs:conj:qu-acs}; this is the
Chapter~\ref{v4-arith:3d-acs-E1bar} residual, inherited here without
modification. The Wilson-Frobenius identity is unconditional at each
finite level; the continuous-group limit is the open problem.
\textbf{Inherited residual.}

\item \emph{Archimedean Wilson-surface realisation.}
\Cref{v4-arith:w9-r1:cor:global-L} requires identification of the
archimedean $\Gamma$-factor with a Wilson-surface expectation at
$\{\infty\}$ via the Arakelov area form
(\Cref{v4-arith:w8-a:rem:wilson-infinity}). Scope: analytic continuation
of the archimedean path integral along Arakelov area flow; tracked in
Chapter~\ref{v4-arith:w8-a:deninger-chiral-homology}.
\textbf{Inherited residual.}

\item \emph{$n\geq 3$ generalisation.}
\Cref{v4-arith:w9-r1:thm:main} is stated for $GL_{2}$, matching the
gauge-group choice of Chapter~\ref{v4-arith:3d-acs-E1bar}. Extension to
$GL_{n}$ goes through
Chapter~\ref{v4-arith:all-genus-koszul}'s genus-$g$ extension with
$G=GL_{2g}(\widehat{\mathbb Z})$; the Morishita functor of
\Cref{v4-arith:w9-r1:thm:morishita-functor} generalises directly to
$G=GL_{n}(\widehat{\mathbb Z})$ by replacing the trace
$\operatorname{tr}$ with the appropriate class function. The
renormalisation residual of item (2) is inherited.
\textbf{Scoped as extension; proof identical modulo replacement of
$GL_{2}$-trace by $GL_{n}$-class function.}
\end{enumerate}

\begin{remark}[the abelian case is closed unconditionally]
\label{v4-arith:w9-r1:rem:abelian-closed}
Collecting: on the abelian-finite-group unramified locus,
\Cref{v4-arith:w9-r1:thm:classical-wilson-frob} proves
$W^{\mathrm{cl}}_{\chi}(K_{p})=\chi(\mathrm{Frob}_{p})$
\emph{unconditionally}. Under
\Cref{v4-arith:3d-acs:conj:qu-acs}, the identity extends to the
continuous-group $GL_{n}(\widehat{\mathbb Z})$ and produces the
finite-prime Frobenius-trace input used by the Deninger skeleton. It
does not construct Deninger's global Frobenius flow on the middle
cohomology. The chapter discharges the finite-level unramified Wilson
calculation and leaves the continuous-group axiom
\Cref{v4-arith:3d-acs:conj:qu-acs}(3) as part of the profinite ACS
renormalisation package.
\end{remark}

\section{Summary}
\label{v4-arith:w9-r1:summary}

\begin{enumerate}
\item The prime knot $K_{p}\subset\overline{\mathrm{Spec}\,\mathbb Z}$
is the Morishita avatar of a codimension-two defect; its meridian is
$I_{p}$ and its longitude is $\mathrm{Frob}_{p}$
(\Cref{v4-arith:w9-r1:def:prime-knot},
\Cref{v4-arith:w9-r1:rem:meridian-longitude}).

\item Classical Kim ACS Wilson loops at finite abelian gauge
group $A$ satisfy, unconditionally,
$\langle W^{\mathrm{cl}}_{\chi}(K_{p})\rangle=\chi(\mathrm{Frob}_{p})$
for unramified $\chi$
(\Cref{v4-arith:w9-r1:thm:classical-wilson-frob}).

\item The Morishita dictionary lifts to a functor
$\Phi_{\mathrm{Mor}}\colon\mathrm{Knot}^{\mathrm{Ar}}\to\mathrm{Conj}(G_{N})$
on the unramified locus; this functor intertwines Wilson insertion
with Frobenius-class evaluation
(\Cref{v4-arith:w9-r1:thm:morishita-functor}).

\item At each finite level $N$ coprime to $p$, the finite-level
Wilson-loop expectation is a finite state sum, with the Frobenius
trace appearing as a class-function insertion
(\Cref{v4-arith:w9-r1:thm:main}(1)), unconditionally.

\item At the profinite limit, the Wilson-loop identity
$W^{\mathrm{qu}}_{\rho}(K_{p})=\operatorname{tr}\rho(\mathrm{Frob}_{p})$
holds conditionally on
\Cref{v4-arith:3d-acs:conj:qu-acs} (the continuous-group
renormalisation), but unconditionally once that conjecture is granted;
no additional open problem appears in the limit step
(\Cref{v4-arith:w9-r1:thm:main}(2)).

\item The local L-factor
$L_{p}(\rho,s)=\det(1-W^{\mathrm{qu}}_{\rho}(K_{p})\,p^{-s})^{-1}$
equals the unramified Satake L-factor, and the global L-function is a
product of Wilson-loop expectations over prime knots plus an
archimedean Wilson-surface factor
(\Cref{v4-arith:w9-r1:thm:local-L-factor},
\Cref{v4-arith:w9-r1:cor:global-L}).

\item Residual obstructions: ramified and wild primes; profinite
renormalisation (inherited); archimedean Wilson-surface realisation
(inherited); $GL_{n}$-generalisation
(\Cref{v4-arith:w9-r1:sec:residual}).
\end{enumerate}

The chapter proves the unramified finite-level Wilson/Frobenius
calculation and conditionally transports it to the continuous-group
limit. In the Deninger chiral candidate of
Chapter~\ref{v4-arith:w8-a:deninger-chiral-homology}, this supplies
only the finite-prime trace component. The global Deninger flow,
archimedean factor, and zero-side spectral operator remain the
spectral-flow/profinite-renormalisation obligations; this chapter does
not resolve them.

\section*{Bibliography of this chapter}
\label{v4-arith:w9-r1:biblio}

The named primary sources for the arguments above:

\begin{itemize}
\item Mazur 1973, \emph{Notes on etale cohomology of number fields},
Ann.~Sci.~ENS~6, for the cohomological 3-dimensionality of
$\mathrm{Spec}\,\mathcal O_{K}$ and the Artin--Verdier
fundamental class.
\item Morishita 2012, \emph{Knots and Primes}, arXiv:0904.3399,
for the knots-primes dictionary: meridian $=$ inertia; longitude
$=$ Frobenius.
\item Witten 1989, \emph{Quantum field theory and the Jones
polynomial}, Commun.~Math.~Phys.~121, for the path-integral
definition of the Wilson loop and its evaluation as a holonomy
character.
\item Kim 2015, \emph{Arithmetic Chern--Simons theory I},
arXiv:1510.05818, for the classical arithmetic Chern--Simons
action on $\overline{\mathrm{Spec}\,\mathcal O_{K}}$ and its
abelian finite state sum.
\item Chern--Dan--Kim--Park 2018, arXiv:1609.03012, for the
non-abelian extension and the Wilson-defect formalism at finite
gauge group.
\item Borger 2009, arXiv:0906.3146, \S8, for the
Bost--Connes Adams $\psi_{p}$ as Frobenius trace on the
crystalline Galois representation.
\item Tate 1979, \emph{Number theory past and present},
Corvallis~II, \S3, for the unramified local Langlands
correspondence at $GL_{1}$.
\item Borel 1976, \emph{Automorphic L-functions}, Corvallis~II,
\S3, and Satake 1963 \emph{Publ.~IHES}~18, for the Satake
parameters and the unramified local L-factor at $GL_{n}$.
\item Laumon 1987, \emph{Transformation de Fourier, constantes
d'equations fonctionnelles et conjecture de Weil},
Publ.~IHES~65, for the epsilon-factor machine at ramified
primes.
\item Fargues--Scholze 2021, arXiv:2102.13459, for the local
geometrisation of the Langlands correspondence (not used in the
proof; referenced for disjoint-source verification).
\end{itemize}
